
\documentclass[b5j,7.5pt,twoside,twocolumn]{tetsujsarticle}
\usepackage{H1授業プリント}
\usepackage{数学ポイント集}
\begin{document}

\setcounter{回番号}{1}
\renewcommand{\回タイトル}{極限⑴}
\twocolumn[
\begin{boxnote}
\hfil \mg{\Large 校内模試に向けて}
\end{boxnote}
\br1]

\hyoudai{はじめに}
本日は校内模試前最終授業です。この半年間で得たものをフルパワーで今回の校内模試にぶつけてください。皆さんが半年に一度のこの校内模試で最大限の力を発揮できるように，少しばかりアドバイスを贈りたいと思います。

\br{1}
\hyoudai{絶対に守るべきこと}

毎週授業補助プリントで告知していますが，数学の校内模試の日時は
\bgt{8/\ihspace{1zw}(\ihspace{1zw})  \ihspace{1zw}:\ihspace{1zw}開始}
です。受験票を確認して，受験番号と教室は覚えたうえで会場にいきましょう。
あとは，携帯以外の時間が分かる何かは必ず持っていってください。皆さんご存知の通り，時計は教室の後ろにしかありません。しかし，試験中に後ろをチラチラ見ていたら試験監督に怪しまれます。ですので，\強調{時間管理のために必ず時計は持参しましょう。}

あとは風邪気味の人はマスクやティッシュ，薬など，各種必需品を持ってくるとよいでしょう。また，シャープペンシルや鉛筆，消しゴムは多めに持っていきましょう。壊れたり落としたりしたときに動揺しないためです。


\br{1}

\hyoudai{試験が始まる前の準備}

解答用紙や問題用紙はそれよりも前から配布されるので，時間に余裕を持って来ましょう。
\bgt{早めに教室に来ることをオススメします。}


やはり，教室に入ってからすぐに試験が始まってしまうと心の整理がつかないまま試験が始まってしまうおそれがあります。\下線{早めに教室に入り，基本事項などをもう一度確認しておいたり，計算テストを直前まで解き直したりしてエンジン全開の状態でテストを迎えましょう。}私オリジナルの計算ドリルテストをやるのもオススメです。

事前に甘いものを口に入れておくと頭が働きやすいですよ。また，テスト前に一度お手洗いにっておきましょう。試験が始まると2時間程度教室から出られなくなるので，テスト前は混む前に早めにお手洗いは済ませておくのがオススメです。

また，解答用紙が配られたら試験開始前に名前とクラスと受験番号を書くことができます。（問題用紙は始めの合図があるまで開かないでください。）%%%%%%%%%%%%Eクラスのところ変える


\br{1}

\hyoudai{試験が始まったら}

試験が始まったら早く解きたい気持ちは非常によく分かります。しかし，実は
\bgt{テストで一番ミスが発生しやすいのは試験開始直後です。}
非常に大切な点なのですが，試験開始直後はみんな必死になって解き始めるのですが，校内模試という非日常な雰囲気のなかで初めからスタートダッシュを必死に決めてしまうと，かえって焦りを生んでしまい逆効果です。

試験が始まったらまずは，ゆっくり鉛筆orシャープペンシルを手にとって，「周りの人はめっちゃ焦ってといているなぁ〜〜」くらいの心持ちで，全ての問題にざっと目を通してから解き始めるくらいが丁度良いのです。最初こそ落ち着いて解き始めましょう。

%\hyoudai{試験中，緊張してしまったら}
%
%
%「緊張」は多くの生徒にとってまさに大敵。試験での緊張は「テスト不安」という言葉にもなっています。「テスト不安」は，試験に失敗することの恐れ，良い点数を取らなければというプレッシャー，友達や親など人からの評価を気にする気持ち，などによって引き起こされる不安のことです。集中力が低下し、普段の力が出せない原因にもなっています。
%
%しかし、駒澤大学文学部心理学科の有光興記教授によると，
%
%\begin{center}
%\強調{適度な緊張状態であればパフォーマンスは良くなる}
%\end{center}
%
%といいます。緊張の高まりとパフォーマンスの出来をグラフに表すと図のような曲線となります。
%
%
%\img{校内模試対策.pdf}
%
%このグラフからも明らかなように，適度な緊張はむしろ試験で最大限のパフォーマンスを発揮する上では非常に良いものとなります。緊張はしすぎてもダメですが，緊張しなさすぎもダメなのです。ですので，\下線{緊張しているからといってヤバい！\hspace{1zw}と焦るのではなく，むしろ「自分はいい状態なんだ」と言い聞かせることで自然と「Bestな状態」にもっていってください。}（勝手に「Bestな状態」になっていきます。）
%
%どうしても過緊張の状態から抜け出せない人は，一度手を止めて深呼吸して目を上下に動かしたりして心を落ち着けましょう。周りが焦って解いている状況を少し見渡してから，「よしやるぞ！」と再開してください。
%
%
\br{1}
\hyoudai{試験時間が残り少なくなったら}

試験時間が残り20分を切ったら，難しい問題にアタックし続けるのではなく，第１問や第２問，第３問といった\強調{小問集合の問題の見直しに時間を充ててください。これは10分〜15分を切ったら全員絶対行ってください。}

冬の校内模試でも皆さんのほとんどの人が感じているでしょうが，最初の小問集合で何十点も落として悔やしんだ人がいたはずです。そしてその人たちの中のほとんどが，見直しに時間を割くことができなかったと言っていました。

解いている途中で投げ出すのは嫌だと思いますが，残り10分を切ったら \強調{どんなに第５問や第６問の問題があと少しで解けそうだとしても，必ず全員第１問〜第３問の小問集合を見直してください。}これは絶対です。校内模試の答案を返却する際，最初の小問集合で大量に計算ミスをして「最後に解き直す時間を取れなかった」という言い訳は許さないのでそのつもりで。

\end{document}

%%%%%%%%%%%%%%%%%%%%%%%%%%%%%%%%%%%%%%%%%%%%%%%%%%%%%%%%%%%%

%\twocolumn[
%\begin{boxnote}
%\hfil \mg{\Large 連絡事項}
%\end{boxnote}
%\br1]
%
%%\標題{連絡事項}
%%
%%\begin{箇条書き}
%%\item あああ
%%\end{箇条書き}
%%
%%
%%\br1
%
%\hyoudai{宿題}
%\msp
%
%\begin{枠囲み}{宿題}
%\begin{enumerate}[\Kaku{enumi}]
%\item \gt{例題の復習（第１回）}
%\item \gt{計算テストの予習}（満点が取れるなら見るだけでOK）
%\item \gt{問題集A・B・C問題}
%\item \gt{問題集の間違え直し}
%\item \gt{計算ドリル}（任意）
%\end{enumerate}
%\end{枠囲み}
%
%\begin{箇条書き}
%\item \bgt{三角関数，指数対数の公式があやふやな人は３ページ以降の付録を読み込み，計算ドリルに取り組むこと}。
%\item 宿題チェックがしやすいようにご協力お願いします。
%\begin{enumerate}[-]
%\item 間違え直しのページには大きく「間違え直し」と記入する。
%\item 宿題の最後のページには「例題◯，問題集◯，間違え直し◯」などとやったことを記入する。
%\end{enumerate}
%\item 次週の居残りの基準。
%\begin{enumerate}[-]
%\item 計算テスト100点未満（チェックなし）
%\item 復習テスト90点以上100点未満（チェックなし）
%\item 復習テスト90点未満（チェックあり）
%\end{enumerate}
%\item \強調{21:00には帰すので，校内模試を理由に休まないでください}。コロナを理由に休むのは仕方がないでしょう。\強調{この授業が配信されますが，復習テストを受けるのと受けないのでは大違いです}。
%\end{箇条書き}
%
%%%%%%%%%%%%%%%%%%%%%%%%%%%%%%%%%%%%%%%%%%%%%%%%%%%%%%%%%%%%%
%
%\twocolumn[
%\begin{columnbox6}{第２回計算テスト追試}
%\mgfamily
%
%目標制限時間：{\Large\強調{4分00秒}}\ihspace{1zw}\→次回は第１回追試と合わせて, {\Large\強調{8分00秒}}で実施。
%\sp
%{\large 「計算テストを疎かにする者, 校内模試の小問集合で泣くべし」}
%\br{-0.3}
%\end{columnbox6}
%\spp]
%
%\img*[c]{re2-1.pdf}
%
%\newcolumn
%
%\img*[c]{re2-2.pdf}
%
%%%%%%%%%%%%%%%%%%%%%%%%%%%%%%%%%%%%%%%%%%%%%%%%%%%%%%%%%%%%%
%
\twocolumn[
\begin{boxnote}
\hfil \mg{\Large 三角関数，指数・対数の公式, 数列}
\end{boxnote}
\br1]

\小見出し{\ajKaku1\ihspace{.5zw}三角関数}

\標題{三角関数と回転角}

\講評段落
ここでは，角が$\theta$のときと$\theta+2n\pi,\,\theta+\pi$などのときとで三角関数の値にどのような関係があるかを説明する。

\ポイント{単位円を描いて，導けるようにすること}

\img[c]{J2-2nd-text-17-02.pdf}

\sp
\講評段落
角$\theta+2n\pi\ (n\text{は整数})$を表す動径は角$\theta$を表す動径と一致する。したがって
\[\text{\ajMaru{1}}
\left\{\begin{array}{l}
\sin(\theta+2n\pi)=\sin\theta\\[5pt]
\cos(\theta+2n\pi)=\cos\theta\\[5pt]
\tan(\theta+2n\pi)=\tan\theta
\end{array}\right.\]


\img[c]{J2-2nd-text-17-03.pdf}

\sp
\講評段落
角$-\theta$を表す動径は角$\theta$を表す動径と
$x$軸に関して対称であるから
\[\text{\ajMaru{2}}
\left\{\begin{array}{l}
\sin(-\theta)=-\sin\theta\\[5pt]
\cos(-\theta)=\cos\theta\\[5pt]
\tan(-\theta)=-\tan\theta
\end{array}\right.\]

\img[c]{J2-2nd-text-17-04.pdf}

\newcolumn
\講評段落
角$\theta+\pi$を表す動径は角$\theta$を表す動径と原点に関して対称であるから
\[\text{\ajMaru{3}}
\left\{\begin{array}{l}
\sin(\theta+\pi)=-\sin\theta\\[5pt]
\cos(\theta+\pi)=-\cos\theta\\[5pt]
\tan(\theta+\pi)=\tan\theta
\end{array}\right.\]

\img[c]{J2-2nd-text-17-05.pdf}

\sp
\講評段落
角$\pi-\theta$を表す動径は角$\theta$を表す動径と$y$軸に関して対称であるから
\[\text{\ajMaru{4}}
\left\{\begin{array}{l}
\sin(\pi-\theta)=\sin\theta\\[5pt]
\cos(\pi-\theta)=-\cos\theta\\[5pt]
\tan(\pi-\theta)=-\tan\theta
\end{array}\right.\]

\img[c]{J2-2nd-text-17-06.pdf}

\newcolumn
\講評段落
角$\theta$を表す表す動径OP上の$\text{P}(x,\,y)$に対して，全体を$\bunsuu{\pi}{2}$回転させると，Pは動径$\text{OP}'$の点$\text{P}'(-y,\,x)$に移る。
よって
\[\text{\ajMaru{5}}
\left\{\begin{array}{l}
\sin\left(\theta+\bunsuu{\pi}{2}\right)=\cos\theta\\[10pt]
\cos\left(\theta+\bunsuu{\pi}{2}\right)=-\sin\theta\\[10pt]
\tan\left(\theta+\bunsuu{\pi}{2}\right)=-\bunsuu{1}{\tan\theta}
\end{array}\right.\]

\img[c]{J2-2nd-text-17-07.pdf}

\sp
\講評段落
角$\theta$を表す表す動径OP上の$\text{P}(x,\,y)$に対して，全体を直線$y=x$に関して対称移動すると，Pは動径$\text{OP}'$の点$\text{P}'(y,\,x)$に移る。
よって
\[\text{\ajMaru{6}}
\left\{\begin{array}{l}
\sin\left(\bunsuu{\pi}{2}-\theta\right)=\cos\theta\\[10pt]
\cos\left(\bunsuu{\pi}{2}-\theta\right)=\sin\theta\\[10pt]
\tan\left(\bunsuu{\pi}{2}-\theta\right)=\bunsuu{1}{\tan\theta}
\end{array}\right.\]

\br1

\標題{加法定理}

\講評段落
三角関数については，２つの角の和や差に関する次の公式が重要である。

\ポイント{加法定理は絶対に覚えること！

これを覚えていないと，他の公式を導けない。}


\begin{枠囲み}{}
\begin{align*}
&\bs{\sin(\alpha\pm\beta)=\sin\alpha\cos\beta\pm\cos\alpha\sin\beta}\\[3pt]
&\bs{\cos(\alpha\pm\beta)=\cos\alpha\cos\beta\mp\sin\alpha\sin\beta}\\[3pt]
&\bs{\tan(\alpha\pm\beta)=\bun{\tan\alpha\pm\tan\beta}{1\mp\tan\alpha\cdot\tan\beta}}\text{（すべて複号同順）}
\end{align*}

\end{枠囲み}


\ポイント{加法定理は余弦定理からも導けるということを覚えておこう。

（20年前の東大入試で出題された…）}

\br{.5}

\noindent
（証明）

\br{-1}
\img<.6>[c]{J3-1st-text-07-01.pdf}

２点$\mathrm{P}\:(\cos\alpha,\,\sin\alpha)$，$\mathrm{Q}\:(\cos\beta,\,\sin\beta)$を考える。このとき，
\begin{align*}
\mathrm{PQ}^2&=(\cos\alpha-\cos\beta)^2+(\sin\alpha-\sin\beta)^2\\
&=2-2(\cos\alpha\cos\beta+\sin\alpha\sin\beta)\text{　……\ajMaru{1}}
\end{align*}

また，$\rmP,\,\rmQ$を原点を中心に反時計回りに$-\beta$だけ回転した点をそれぞれ$\rmP',\,\rmQ'$とすると，\\$\rmP'\left(\cos(\alpha-\beta),\,\sin(\alpha-\beta)\right),\ \rmQ'(1,\,0)$となるので，
\begin{align*}
\mathrm{P'Q'}^2&=\left\{\cos(\alpha-\beta)-1\right\}^2+\sin^2(\alpha-\beta)\\
&=2-2\cos(\alpha-\beta)\text{　……\ajMaru{2}}
\end{align*}

$\caprm PQ=P'Q'$なので，\ajMaru{1}\ajMaru{2}より

\begin{枠囲み}{}
\[\bs{\cos(\alpha-\beta)=\cos\alpha\cos\beta+\sin\alpha\sin\beta}\text{……\ajMaru{3}}\]
\end{枠囲み}

\ajMaru{3}の$\beta$を$-\beta$にすると
\[\cos(\alpha+\beta)=\cos\alpha\cos(-\beta)+\sin\alpha\sin(-\beta)\]

$\cos(-\beta)=\cos\beta$，$\sin(-\beta)=-\sin\beta$より

\begin{枠囲み}{}
\[\bs{\cos(\alpha+\beta)=\cos\alpha\cos\beta-\sin\alpha\sin\beta}\]
\end{枠囲み}

\ajMaru{3}の$\alpha$を$\bun{\pi}{2}-\alpha$にすると，$\cos\left(\bun{\pi}{2}-\alpha-\beta\right)=\sin(\alpha+\beta)$より，
\br{0.2}
\[\sin(\alpha+\beta)=\cos\left(\bun{\pi}{2}-\alpha\right)\cos\beta+\sin\left(\bun{\pi}{2}-\alpha\right)\sin\beta\]

\br{0.2}
$\cos\left(\bun{\pi}{2}-\alpha\right)=\sin\alpha$，$\sin\left(\bun{\pi}{2}-\alpha\right)=\cos\alpha$より

\begin{枠囲み}{}
\[\bs{\sin(\alpha+\beta)=\sin\alpha\cos\beta+\cos\alpha\sin\beta}\]
\end{枠囲み}

$\beta$を$-\beta$にして$\cos(-\beta)=\cos\beta$，$\sin(-\beta)=-\sin\beta$を用いると

\begin{枠囲み}{}
\[\bs{\sin(\alpha-\beta)=\sin\alpha\cos\beta-\cos\alpha\sin\beta}\]
\end{枠囲み}

また，正接に関しては，以下複号同順として，
\br{0.4}
\[\tan(\alpha\pm\beta)=\bun{\sin(\alpha\pm\beta)}{\cos(\alpha\pm\beta)}=\bun{\sin\alpha\cos\beta\pm\cos\alpha\sin\beta}{\cos\alpha\cos\beta\mp\sin\alpha\sin\beta}\]

\br{0.4}
右辺の分母分子を$\cos\alpha\cos\beta$で割って
\br{0.4}
\[\tan(\alpha\pm\beta)=\bun{\bun{\sin\alpha\cos\beta}{\cos\alpha\cos\beta}\pm\bun{\cos\alpha\sin\beta}{\cos\alpha\cos\beta}}{1\mp\bun{\sin\alpha\sin\beta}{\cos\alpha\cos\beta}}=\bun{\bun{\sin\alpha}{\cos\alpha}\pm\bun{\sin\beta}{\cos\beta}}{1\mp\bun{\sin\alpha}{\cos\alpha}\cdot\bun{\sin\beta}{\cos\beta}}\]

\begin{枠囲み}{}
\[\bs{\tan(\alpha\pm\beta)=\bun{\tan\alpha\pm\tan\beta}{1\mp\tan\alpha\cdot\tan\beta}}\text{（複号同順）}\]
\end{枠囲み}

\newcolumn

\br1

\標題{倍角公式}


\ポイント{登場頻度が高いので覚えるべき。忘れてしまったら, 

加法定理において$\;\beta=\alpha\;$とおいて導く。}


\begin{枠囲み}{倍角公式}
\[\sin2\alpha&=2\sin\alpha\cos\alpha\\
\cos2\alpha&=\cos^2\alpha-\sin^2\alpha\\
&=2\cos^2\alpha-1\quad(\because\sin^2\alpha+\cos^2\alpha=1)\\
&=1-2\sin^2\alpha\quad(\because\sin^2\alpha+\cos^2\alpha=1)\\
\tan2\alpha&=\bun{2\tan\alpha}{1-\tan^2\alpha}\]
\end{枠囲み}


\begin{枠囲み}{補充問題}
$\cos\theta=\bun{4}{5}\;\left(0<\theta<\bun{\pi}{2}\right)$のとき，$\sin2\theta,\,\cos2\theta,\,\tan2\theta$を求めよ。
\br{0.2}
\end{枠囲み}

\見出し{解答}

$\cos\theta=\bun45$より，
\[\sin^2\theta=1-\cos^2\theta=1-\left(\bun45\right)^2=\bun9{25}\]

$0<\theta<\bun{\pi}{2}$であるから，
\[\sin\theta>0\]

よって，
\[\sin\theta=\bun35\]

ここで，倍角公式より，
\begin{align*}
\sin2\theta&=2\sin\theta\cos\theta\\
&=2\cdot\bun35\cdot\bun45=\bs{\bun{24}{25}}\kotae \\\cos2\theta&=\cos^2\theta-\sin^2\theta\\
&=\left(\bun45\right)^2-\left(\bun35\right)^2=\bs{\bun7{25}}\kotae\\
\tan2\theta&=\bun{\sin2\theta}{\cos2\theta}=\bs{\bun{24}{7}}\kotae
\end{align*}


\br{0.5}
\begin{枠囲み}{倍角公式の応用２}%%%%%%%%%%%%%%%%%%%%%%%%%%%%%%%%%%%%%　枠囲み06（J3-08枠囲み11）
$0\leqq x<2\pi$のとき，次の方程式・不等式を解け。
\br{0.2}
\設問カッコ{1}{$\sin2x=\sqrt{3}\sin x$}
\br{0.2}
\設問カッコ{2}{$\cos2x-\sin x>1$}
\end{枠囲み}
\br{0.2}
\begin{指針}
種類・角とも揃っていない。そこで，%\gt{【枠囲み10】}と同様，


\解答小問{
\shoumonKakko{1}%%%%%%%%%%%%%%%%%%%%(1)
$\sin2x=2\sin x\cos x$より，
\begin{align*}
\text{（与式）}&\Leftrightarrow 2\sin x\cos x=\sqrt3\sin x \\
&\Leftrightarrow \sin x\rootbracketl(2\cos x-\sqrt3\,\rootbracketr)=0\\
&\Leftrightarrow \sin x=0,\ \cos x=\bun{\sqrt3}2
\end{align*}

$0\leqq x<2\pi$であるから，
\[\bs{x=0,\ \pi,\ \bun{\pi}{6},\ \bun{11\pi}{6}}\kotae\]

\shoumonKakko{2}%%%%%%%%%%%%%%%%%%%%(2)
$\cos2x=1-2\sin^2x$より，
\begin{align*}
\text{（与式）}&\Leftrightarrow 1-2\sin^2x-\sin x>1\\
&\Leftrightarrow 2\sin x\left(\sin x+\bun12\right)<0\\
&\Leftrightarrow -\bun12<\sin x<0
\end{align*}


$0\leqq x<2\pi$であるから，
\[\bs{\pi<x<\bun{7\pi}{6},\ \bun{11\pi}{6}<x<2\pi}\kotae\]
}



\begin{key}
\br{0.2}
まず，種類・角を統一する
\br{0.2}
\end{key}

\noindent
ことを考える。

\end{指針}

\br1

\標題{半角公式}

\ポイント{登場頻度が高いので覚えるべき。忘れてしまったら，

$\cos$の倍角公式の$\alpha$を$\bun{\alpha}{2}$とおいて導く。}

\begin{枠囲み}{半角公式}
\[&\sin^2\bun{\alpha}{2}\hspace{2.5pt}=\:\bun{1-\cos\alpha}{2}\\
&\cos^2\bun{\alpha}{2}\hspace{1.5pt}=\:\bun{1+\cos\alpha}{2}\\
&\tan^2\bun{\alpha}{2}=\:\bun{1-\cos\alpha}{1+\cos\alpha}\]
\end{枠囲み}


\br{0.5}
\begin{枠囲み}{半角公式}%%%%%%%%%%%%%%%%%%%%%%%%%%%%%%%%%%%%%%%%%%%%　枠囲み07（J3-08枠囲み12）
\br{0.2}
$\sin\theta=-\bun35\;\left(\pi<\theta<\bun{3}{2}\pi\right)$のとき，$\sin\bun{\theta}{2}$，$\cos\bun{\theta}{2}$を求めよ。
\br{0.2}
\end{枠囲み}
\br{0.5}
\begin{指針}
エースの式から$\cos\theta$を求め，半角公式に代入，$\theta$の範囲から，$\sin\bun{\theta}{2}$，$\cos\bun{\theta}{2}$の符号を確定できることにも注意。
\end{指針}

$\sin\theta=-\bun35$より，
\[\cos^2\theta=1-\sin^2\theta=1-\left(-\bun35\right)^2=\bun{16}{25}\]

$\pi<\theta<\bun32\pi$であるから，
\[\cos\theta<0\]

よって，
\[\cos\theta=-\bun45\]

半角公式より，
\begin{align*}
\sin^2\bun{\theta}{2}=\bun{1-\cos\theta}2=\bun{1-\left(-\bun45\right)}2=\bun9{10}\\
\cos^2\bun{\theta}{2}=\bun{1+\cos\theta}2=\bun{1+\left(-\bun45\right)}2=\bun1{10}\\
\end{align*}

ここで，$\pi<\theta<\bun32\pi$であるから，
\[\bun{\pi}{2}<\bun{\theta}2<\bun34\pi\]

よって，
\[\cos\bun{\theta}2<0,\ \sin\bun{\theta}2>0\]

以上から，
\[\bs{\sin\bun{\theta}2=\bun{3\sqrt{10}}{10},\ \cos\bun{\theta}2=-\bun{\sqrt{10}}{10}}\kotae\]

\br1
\noindent

\br1

\標題{３倍角の公式}

\ポイント{【参考】の語呂合わせで覚えた。忘れてしまったら，

加法定理の$\beta$を$2\alpha$とおいて導く。}

%\begin{枠囲み}{補充問題１}%%%%%%%%%%%%%%%%%%%%%%%%%%%%%%%%%%%%%%%%%　枠囲み08（J3-08枠囲み13）
%
%$\sin3\alpha=3\sin\alpha-4\sin^3\alpha$であることを証明せよ。
%
%\end{枠囲み}


\begin{枠囲み}{３倍角の公式}
\begin{align*}
&\sin3\alpha=3\sin\alpha-4\sin^3\alpha\\[2pt]
&\cos3\alpha=4\cos^3\alpha-3\cos\alpha
\end{align*}
\end{枠囲み}

\begin{答案段落}{導出}
加法定理より，
\begin{align*}
\sin3\alpha&=\sin(2\alpha+\alpha)\\
&=\sin2\alpha\cos\alpha+\cos2\alpha\sin\alpha\\
&=2\sin\alpha\cos^2\alpha+(1-2\sin^2\alpha)\sin\alpha\\
&=2\sin\alpha(1-\sin^2\alpha)+\sin\alpha-2\sin^3\alpha\\
&=3\sin\alpha-4\sin^3\alpha
\end{align*}

よって，
\[\sin3\alpha=3\sin\alpha-4\sin^3\alpha\qed\]
\end{答案段落}

\begin{枠囲み}{参考：語呂合わせ}
\hspace{4zw}$\sin3\alpha\quad=\quad\underline{3}\:\underline{\sin\alpha}\qquad\underline{-}\qquad\underline{4}\:\underline{\sin^3\alpha}$

\hspace{9.5zw}{\scriptsize \underline{サン}\:\underline{シャイン}\;\;\;\;\underline{落ちて}\;\;\;\;\underline{よ}\hspace{-0.5pt}かぜが\;\underline{さ}\hspace{-0.5pt}びしいな}

\sp
\qquad$\cos3\alpha$は$\sin3\alpha$において$\sin\alpha$を$\cos\alpha$に変え全体にマイナスをつければよい。
\end{枠囲み}

\newcolumn

\標題{和積の公式}

\ポイント{【参考】の語呂合わせ（もはや呪文）で覚えた。

加法定理を辺々足し引きして導く方法も覚えるべき。

和積の公式は微分で使います！}%注：$f'(x)$を積の形にするとき

\begin{枠囲み}{和積の公式}
\begin{align*}
&\sin \alpha+\sin \beta=2\sin \bun{\alpha+\beta}2\cos \bun{\alpha-\beta}2\\[2pt]
&\sin \alpha-\sin \beta=2\cos \bun{\alpha+\beta}2\sin \bun{\alpha-\beta}2\\[2pt]
&\cos \alpha+\cos \beta=2\cos \bun{\alpha+\beta}2\cos \bun{\alpha-\beta}2\\[2pt]
&\cos \alpha-\cos \beta=-2\sin \bun{\alpha+\beta}2\sin \bun{\alpha-\beta}2
\end{align*}
\end{枠囲み}

\sp
\講評段落
加法定理を辺々足し引きすると導くことができます。

\begin{答案段落}{導出}
\[&\sin (\alpha+\beta)=\sin \alpha\cos \beta+\cos \alpha\sin \beta\\
&\sin (\alpha-\beta)=\sin \alpha\cos \beta-\cos \alpha\sin \beta\]

辺々足して, 
\[\sin (\alpha+\beta)+\sin (\alpha-\beta)=2\sin \alpha\cos \beta\]

辺々引いて, 
\[\sin (\alpha+\beta)-\sin (\alpha-\beta)=2\cos \alpha\sin \beta\]

$\alpha+\beta$を$x$, $\alpha-\beta$を$y$に置き換えると, $\alpha=\bun{x+y}2$, $\beta=\bun{x-y}2$と表せるので, 

\[&\sin x+\sin y=2\sin \bun{x+y}2\cos \bun{x-y}2\\[2pt]
&\sin x-\sin y=2\cos \bun{x+y}2\sin \bun{x-y}2\]
が導かれる。

\[&\cos (\alpha+\beta)=\cos \alpha\cos \beta-\sin \alpha\sin \beta\\
&\cos (\alpha-\beta)=\cos \alpha\cos \beta+\sin \alpha\sin\beta\]

辺々足して, 
\[\cos (\alpha+\beta)+\cos (\alpha-\beta)=2\cos \alpha\cos \beta\]

辺々引いて, 
\[\cos (\alpha+\beta)-\cos (\alpha-\beta)=-2\sin \alpha\sin \beta\]

$\alpha+\beta$を$x$, $\alpha-\beta$を$y$に置き換えると, $\alpha=\bun{x+y}2$, $\beta=\bun{x-y}2$と表せるので, 

\[&\cos x+\cos y=2\cos \bun{x+y}2\cos \bun{x-y}2\\
&\cos x-\cos y=-2\sin \bun{x+y}2\sin \bun{x-y}2\]
が導かれる。
\end{答案段落}

\begin{枠囲み}{参考：語呂合わせ}
\hspace{4zw}$\underline{\sin\alpha}\quad\underline{+}\quad\underline{\sin\beta}\quad\underline{=}\quad\underline{2}\quad\underline{\sin\bun{\alpha+\beta}2}\quad\underline{\cos\bun{\alpha-\beta}2}$
\br{.3}
\hspace{4zw}{\scriptsize \;\underline{シン}\hspace{2zw}\underline{足す}\hspace{2zw}\underline{シン}\hspace{2zw}\underline{は}\quad\underline{に}\hspace{4zw}\underline{しん}\hspace{3zw}の\hspace{3zw}\underline{子}}
\sp
同様に, 「シン引くシンはにコスシン」「コス足すコスはにコスコス」「コス引くコスはマイナスにシンシン」と覚える（もはや呪文）。
\end{枠囲み}

\newcolumn

\標題{積和の公式}

\ポイント{和積の公式を右辺から左辺にたどれば導くことができる。

積分で頻繁に使うので, 使ううちに覚えていく。}

\begin{枠囲み}{積和の公式}
\begin{align*}
&\sin \alpha\cos \beta =\bun12\left\{\sin (\alpha+\beta)+\sin (\alpha-\beta)\right\}\\[2pt]
&\cos \alpha\sin \beta =\bun12\left\{\sin (\alpha+\beta)-\sin (\alpha-\beta)\right\}\\[2pt]
&\cos \alpha\cos \beta =\bun12\left\{\cos (\alpha+\beta)+\cos (\alpha-\beta)\right\}\\[2pt]
&\sin \alpha\sin \beta =-\bun12\left\{\cos (\alpha+\beta)-\cos (\alpha-\beta)\right\}
\end{align*}
\end{枠囲み}


\begin{答案段落}{導出}
和積の公式
\[&\sin \alpha+\sin \beta=2\sin \bun{\alpha+\beta}2\cos \bun{\alpha-\beta}2\\
\Leftrightarrow &\sin \bun{\alpha+\beta}2\cos \bun{\alpha-\beta}2=\bun12(\sin \alpha+\sin \beta)\]
において, $\bun{\alpha+\beta}2=x$, $\bun{\alpha-\beta}2=y$とすると, $\alpha=x+y$, $\beta=x-y$なので, 

\[\sin x\cos y=\bun12\{\sin (x+y)+\sin (x-y)\}\]
が導かれる。
\spp
もう一つやっておこう。
\[&\cos \alpha-\cos \beta=-2\sin \bun{\alpha+\beta}2\sin \bun{\alpha-\beta}2\\
\Leftrightarrow &\sin \bun{\alpha+\beta}2\sin \bun{\alpha-\beta}2=-\bun12(\cos \alpha-\cos \beta)\]
において, $\bun{\alpha+\beta}2=x$, $\bun{\alpha-\beta}2=y$とすると, $\alpha=x+y$, $\beta=x-y$なので, 
\[\sin x\sin y=-\bun12\{\cos (x+y)-\cos (x-y)\}\]
が導かれる。
\end{答案段落}

\br1

\標題{}

以上大事なのは

\ポイント{公式を忘れてしまっても, 全て加法定理から導けるということである。

そのため, 加法定理だけは絶対に覚えなければならない。}

ということだ。勿論, 全ての公式を覚えるに越したことはない。特に, 

\ポイント{倍角公式・半角公式・積和の公式

の３つは\bgt{積分計算}で頻繁に使っていく。}

ので, 早いうちに覚えるべきである。

%%%\newpage
%%%\newcolumn


\clearpage 

\小見出し{\ajKaku2\ihspace{.5zw}指数・対数}

\標題{指数の拡張}

\講評段落
$m$，$n$が整数のとき，次のような指数法則が成り立つことは，すでに学んだ。
\begin{align*}
a^ma^n&=a^{m+n}\\
\left(a^m\right)^n&=a^{mn}\\
\left(ab\right)^n&=a^nb^n
\end{align*}

特に，
\begin{align*}
&a^0=1\\
&a^{-n}=\bun{1}{a^n}
\end{align*}
のように定義された。

\sp
\講評段落
ここで，指数を有理数全体まで拡張すると，$a>0$，$m$を任意の整数，$n$を正の整数として，
\[\left(a^{\bun{m}{n}}\right)^n=a^{\bun{m}{n}\times n}=a^m\]

したがって，$a^{\bun{m}{n}}$は$a^m$の$n$乗根の正のもの，すなわち
\[a^{\bun{m}{n}}=\sqrt[n\:]{\,a^m\,}\]
と定めよう。

\sp
\講評段落
このように定めると，指数法則も拡張した形で成立する。

\begin{枠囲み}{}
　$a>0$，$b>0$　　$r$，$s$：有理数
\[a^ra^s=a^{r+s}\hspace{3zw}\left(a^r\right)^s=a^{rs}\hspace{3zw}\left(ab\right)^r=a^rb^r\]
\end{枠囲み}

ただし，
\begin{枠囲み}{}
\begin{align*}
&a^0=1\hspace{3zw}a^{-n}=\bun{1}{a^n}\\[5pt]
&a^{\bun{m}{n}}=\sqrt[n\:]{\,a^m\,}\qquad\text{$m$は任意の整数，$n$は正の整数}
\end{align*}
\end{枠囲み}

\noindent ㊟これらは，指数計算が上手く行くように新たに定義したものなので，証明を求めてはならない。

\br1

\標題{対数とその性質}

\img*<0.8>[C]{J3-1st-text-11-01.pdf}

\sp
\講評段落
指数関数$\,y=a^x\,$のグラフからわかるように，$x$と正の数$y$とは$1$対$1$対応している。つまり正の数$R$に対して$a^r=R$をみたす実数$r$がただ一つ決まる。この$r$を$\log_aR$と書き，$a$を\bgt{底}とする$R$の\bgt{対数}という。

\newcolumn
すなわち，

\begin{枠囲み}{}
\[\gt{$a^r=R\ \Leftrightarrow\ r=\log_aR$}\qquad\text{（$a>0$，$a\neq1$，$R>0$）}\]
\end{枠囲み}

\noindent
と定義する。このとき，$R$を$a$を底とする$r$の\bgt{真数}という。

したがって，真数$R$は\bgt{つねに正}にならねばならず，底$a$は\bgt{１に等しくない正の数}でなければならない。

（$a=1$のとき，$\,y=a^x\,$をみたす$x$と$y$は$1$対$1$対応しない。）
%}

\begin{枠囲み}{対数の基本法則}
底$a$，$b$は$1$ではない正の数，真数$R$，$S$は正の数，$t$は実数
\begin{align*}
&\text{\ajMaru{1}　$\log_a1=0,\quad\log_aa=1$}\\[5pt]
&\text{\ajMaru{2}　$\log_aRS=\log_aR+\log_aS$}\\[5pt]
&\text{\ajMaru{3}　$\log_a\bun{R}{S}=\log_aR-\log_aS$}\\[5pt]
&\text{\ajMaru{4}　$\log_aR^t=t\log_aR$}\\[5pt]
&\text{\ajMaru{5}　$\log_aR=\bun{\log_bR}{\log_ba}$　（底の変換公式）}\\[5pt]
&\text{\ajMaru{6}　$a^{\log_aR}=R$}
\end{align*}
\end{枠囲み}

\見出しぶら下げ{\bgt{（参考）}}{\ajMaru{1}〜\ajMaru{6}の性質は指数と対数の基本関係$\,a^r=R\:\Leftrightarrow\:r=\log_aR\,$を用いて，対数の定義や指数法則から導かれる。}

\br{0.5}
\見出しぶら下げ[2zw]{\ajMaru{1}：}{$\;a^0=1$，$a^1=a$より明らか}

\br{0.5}
\見出しぶら下げ[2zw]{\ajMaru{2}，\ajMaru{3}：}{$a^r=R,\ a^s=S$とすると
\[r=\log_aR,\ s=\log_aS\]

辺々加えて
\[r+s=\log_aR+\log_aS\]

また，指数法則より，$a^{r+s}=a^ra^s$であり，$a^ra^s=RS$であるから，$a^{r+s}=RS$。これより，
\[r+s=\log_aRS\]

以上より，
\[\log_aRS=\log_aR+\log_aS\]

\br{0.5}
同様に，指数法則より，$a^{r-s}=\bun{a^r}{a^s}$であり，$\bun{a^r}{a^s}=\bun{R}S$であるから，$a^{r-s}=\bun{R}S$。これより，
\[r-s=\log_a\bun{R}S\]

ここで，
\[r=\log_aR,\ s=\log_aS\]
を辺々引いて，
\[r-s=\log_aR-\log_aS\]

以上より，
\[\log_a\bun{R}S=\log_aR-\log_aS\]
}

\br{0.5}
\見出しぶら下げ[2zw]{\ajMaru{4}：}{$\;a^r=R$とすると$a^{rt}=\left(a^r\right)^t=R^t$より
\[rt=\log_aR^t\]

$r=\log_aR$なので，
\[\log_aR^t=t\log_aR\]
}

\br{0.5}
\見出しぶら下げ[2zw]{\ajMaru{5}：}{$\log_bR=p$，$\log_ba=q$とおくと
\[R=b^p,\ a=b^q\]

$a^p=(b^q)^p=(b^p)^q=R^q$なので
\[a^{\bun{p}{q}}=R\]

両辺の底$a$の対数をとって，
\[\bun{p}{q}=\log_aR\]

$\log_bR=p$，$\log_ba=q$を代入して，
\[\log_aR=\bun{\log_bR}{\log_ba}\]

\[\left[\begin{array}{l}
特に\;R=b\;のとき\;\log_ab=\bun{\log_bb}{\log_ba}=\bun{1}{\log_ba}\\[7pt]
分母を払って\;\;\underline{\log_ba\log_ab=1}\end{array}\right]\]
}
\br{0.5}
\見出しぶら下げ[2zw]{\ajMaru{6}：}{$a^r=R$とすると両辺の底$a$の対数をとって，
\[r=\log_aR\]

この$r$を$a^r=R$に代入して，
\[a^{\log_aR}=R\]
}

\newpage
\小見出し{\ajKaku3\ihspace{.5zw}数列}
\begin{枠囲み}{和の応用１}%%%%%%%%%%%%%%%%%%%%%%%%%%%%%%%%%%%%%%%%%%%%%%%%%%%　枠囲み12（J3-4枠囲み17）
$\Bun{1}{1\cdot2}+\bunsuu{1}{2\cdot3}+\bunsuu{1}{3\cdot4}+\cdots\cdots\cdots+\bunsuu{1}{n\cdot(n+1)}$　を計算せよ。
\end{枠囲み}

一般項が$\Bun{1}{k\cdot(k+1)}$である数列の初項から第$n$項までの和と考える。

$\bunsuu{1}{k\cdot(k+1)}=\bunsuu{1}{k}-\bunsuu{1}{k+1}$の変形を覚えておこう。

\br{0.3}
\begin{key}
\br{0.2}
\textbf{\textgt{分数で表された数列の和の求め方}}

\br{0.3}
　　　「和の中抜け」が使えるように一般項を変形する。
\end{key}
任意の自然数$k$に対し，
\[
\bun{1}{k(k+1)}=\bun1k-\bun1{k+1}
\]
が成立するので，
\begin{align*}
&\bun1{1\cdot 2}+\bun1{2\cdot 3}+\bun1{3\cdot 4}+\cdots +\bun{1}{n(n+1)}\\
={}&\left(\bun11-\bun12\right)+\left(\bun12-\bun13\right)+\left(\bun13-\bun14\right)+\cdots\\
&\freqn{\cdots +\left(\bun1{n-1}-\bun1n\right)+\left(\bun1n-\bun1{n+1}\right)}\\
={}&1-\bun{1}{n+1}\\
={}&\bs{\bun{n}{n+1}}\kotae
\end{align*}


\begin{枠囲み}{練習問題}
次の和を求めよ。
\一段組小問{⑴\hspace{1zw}$1\cdot2+2\cdot2^{2}+3\cdot2^{3}+4\cdot2^{4}+\cdots+n\cdot2^{n}$}
\一段組小問{⑵\hspace{1zw}$\Bun{11}{1\cdot2\cdot3}+\bunsuu{16}{2\cdot3\cdot4}+\cdots+\bunsuu{5n+6}{n(n+1)(n+2)}$}
\一段組小問{⑶\hspace{1zw}$1\cdot2\cdot3+2\cdot3\cdot4+3\cdot4\cdot5+\cdots+n(n+1)(n+2)$}

\end{枠囲み}

\解答小問{%%%%%%%%%%%%%%%%%%%%%% (1)
\shoumonKakko1求める和を$S$とすると，

$\left\{\begin{array}{lllll}S=1\cdot2+&2\cdot2^2+&3\cdot2^3+\cdots&+n\cdot2^n&\text{　……\ajMaru1}\\
2S=&1\cdot2^2+&2\cdot2^3+\cdots&+(n-1)\cdot2^n&+n\cdot2^{n+1}\\
&&&&\text{　……\ajMaru2}
\end{array}\right.$
$\text{\ajMaru2}-\text{\ajMaru1}$より，
\begin{align*}
S&=-(1\cdot2+1\cdot2^2+1\cdot2^3+\cdots+1\cdot2^n)+n\cdot2^{n+1}\\&=-\bun{2(1-2^n)}{1-2}+n\cdot2^{n+1}\\&=\bs{(n-1)\cdot2^{n+1}+2}
\kotae\end{align*}
}

\解答小問{%%%%%%%%%%%%%%%%%%%%%% (2)
\begin{align*}\alignleftrightarrowKakko2&\bun{11}{1\cdot2\cdot3}+\bun{16}{2\cdot3\cdot4}+\cdots+\bun{5n+6}{n(n+1)(n+2)}\\
={}&\sum_{k=1}^n\bun{5k+6}{k(k+1)(k+2)}\\
={}&\sum_{k=1}^n\left\{\bun3{k(k+1)}+\bun2{(k+1)(k+2)}\right\}\\
={}&3\sum_{k=1}^n\left(\bun1k-\bun1{k+1}\right)+2\sum_{k=1}^n\left(\bun1{k+1}-\bun1{k+2}\right)\\
={}&3\left\{\left(\bun11-\bun12\right)+\left(\bun12-\bun13\right)+\cdots+\left(\bun1n-\bun1{n+1}\right)\right\}\\
&\freqn{+2\left\{\left(\bun12-\bun13\right)+\left(\bun13-\bun14\right)+\cdots+\left(\bun1{n+1}-\bun1{n+2}\right)\right\}}\\
={}&3\left(\bun11-\bun1{n+1}\right)+2\left(\bun12-\bun1{n+2}\right)\\
={}&3-\bun3{n+1}+1-\bun2{n+2}=4-\bun{5n+8}{(n+1)(n+2)}\\
={}&\bs{\bun{n(4n+7)}{(n+1)(n+2)}}\kotae
\end{align*}
}

\解答小問{%%%%%%%%%%%%%%%%%%%%%% (3)
\begin{align*}\alignleftrightarrowKakko3&1\cdot2\cdot3+2\cdot3\cdot4+3\cdot4\cdot5+\cdots+n(n+1)(n+2)\\
={}&\sum_{k=1}^nk(k+1)(k+2)\\
={}&\bun14\sum_{k=1}^n\{k(k+1)(k+2)(k+3)-(k-1)k(k+1)(k+2)\}\\
={}&\bun14[(1\cdot2\cdot3\cdot4-0\cdot1\cdot2\cdot3)\\&+(2\cdot3\cdot4\cdot5-1\cdot2\cdot3\cdot4)+\cdots\\
&+\{n(n+1)(n+2)(n+3)-(n-1)n(n+1)(n+2)\}]\\
={}&\bun14\{n(n+1)(n+2)(n+3)-0\cdot1\cdot2\cdot3\}\\
={}&\bs{\bun14n(n+1)(n+2)(n+3)}\kotae
\end{align*}
}

\見出し{\ajKakko3別解}%\見出し は\解答小問 の外側に書いて下さい。
\解答小問{
\noindent
\begin{align*}
&1\cdot2\cdot3+2\cdot3\cdot4+3\cdot4\cdot5+\cdots+n(n+1)(n+2)\\
={}&\sum_{k=1}^nk(k+1)(k+2)\\
={}&\sum_{k=1}^n(k^3+3k^2+2k)\\
={}&\left\{\bun12n(n+1)\right\}^2+3\cdot\bun{n(n+1)(n+2)}6+2\cdot\bun{n(n+1)}2\\
={}&\bs{\bun14n(n+1)(n+2)(n+3)}\kotae
\end{align*}

}

\begin{枠囲み}{練習問題}
次の数列$\left\{a_{n}\right\}$の一般項を求めよ。
\一段組小問{⑴\hspace{1zw}$a_{1}=1,\ a_{n+1}=a_{n}+2^{n+1}$}
\一段組小問{⑵\hspace{1zw}$a_{1}=1,\ a_{n+1}=a_{n}+\Bun{1}{n(n+1)(n+2)}$}
\一段組小問{⑶\hspace{1zw}$a_1=3,\ a_{n+1}=2a_{n}-2$}

\end{枠囲み}

\解答小問{%%%%%%%%%%%%%%%%%%%%%% (1)
\shoumonKakko1数列$\{a_n\}$の階差数列は$\{2^{n+1}\}$であるので，$n\geqq2$のとき，

\[a_n=1+\sum_{k=1}^{n-1}2^{k+1}=1+\bun{4(1-2^{n-1})}{1-2}=2^{n+1}-3\]

$n=1$のとき，$a_1=1$となるので，$n=1$のときもこれでよい。

よって，

\[a_n=\bs{2^{n+1}-3\ (n\geqq1)}\kotae\]
}

\解答小問{%%%%%%%%%%%%%%%%%%%%%% (2)
\shoumonKakko2数列$\{a_n\}$の階差数列は$\left\{\bun1{n(n+1)(n+2)}\right\}$であるので，$n\geqq2$のとき，
\begin{align*}
a_n&=1+\sum_{k=1}^{n-1}\bun1{k(k+1)(k+2)}\\&=1+\bun12\sum_{k=1}^{n-1}\left\{\bun1{k(k+1)}-\bun1{(k+1)(k+2)}\right\}\\
&=1+\bun12\left[\left(\bun1{1\cdot2}-\bun1{2\cdot3}\right)+\left(\bun1{2\cdot3}-\bun1{3\cdot4}\right)+\right.\\
&\freqn{\left.\cdots+\left\{\bun1{(n-1)n}-\bun1{n(n+1)}\right\}\right]}\\
&=1+\bun12\left\{\bun12-\bun1{n(n+1)}\right\}=\bun{5n^2+5n-2}{4n(n+1)}
\end{align*}

$n=1$のとき，$a_1=1$となるので，$n=1$のときもこれでよい。

よって，
\[a_n=\bs{\bun{5n^2+5n-2}{4n(n+1)}\ (n\geqq1)}\kotae\]
}


\解答小問{%%%%%%%%%%%%%%%%%%%%%% (3)
\[\alignKakko{3}a_{n+1}=2a_n-2\式番号{1}\]

ここで，$\alpha=2\alpha-2$を解くと，$\alpha=2$なので，
\[2=2\cdot 2-2\式番号{2}\]

\ajMaru{1}$-$\ajMaru{2}より，
\[a_{n+1}-2=2(a_n-2)\式番号{3}\]

ここで，$b_n=a_n-2$とおくと，\ajMaru{3}より，
\[b_{n+1}=2b_n\]

$b_1=a_1-2=1$であるから，数列$\left\{b_n\right\}$は，初項$1$，公比$2$の等比数列をなす。よって，
\begin{align*}
&b_n=1\cdot2^{n-1}\\
\LLeftrightarrow &a_n-2=2^{n-1}\\
\LLeftrightarrow &\bs{a_n=2^{n-1}+2\ (n\geqq1)}\kotae
\end{align*}
}




\section{等比数列}%%%%%%%%%%%%%%%%%%%%%%%%%%%%%%%%%%%%%%%%%%%%%%%%%%%%%%%%%%　§３
\subsection{等比数列}%%%%%%%%%%%%%%%%%%%%%%%%%%%%%%%%%%%%%%%%%%%%%%%%%%%%%%%　3-1
\[\text{$1$，$2$，$4$，$8$，$16$，……………}\]
という数列は，初項$1$に次々に一定の数$2$をかけて作られる。

このようにすべての自然数$n$に対して\;$a_{n+1}=ra_n$（$r$は$0$でない定数）となる数列を\textbf{\textgt{等比数列}}という。$r$を等比数列の\textbf{\textgt{公比}}という。

上の数列で考えると，
\br{0.5}
\[\begin{array}{ll}
第1項&\hspace{3zw}1\:\cdot\:2^0\\[5pt]
第2項&\hspace{3zw}1\:\cdot\:2^1\\[5pt]
第3項&\hspace{3zw}1\:\cdot\:2^2\\[-2pt]
\makebox[3zw][c]{$\vdots$}&\hspace{4zw}\vdots\\[-7pt]
\makebox[3zw][c]{$\vdots$}&\hspace{4zw}\vdots\\[-7pt]
\makebox[3zw][c]{$\vdots$}&\hspace{4zw}\vdots\\
第n項&\hspace{3zw}1\:\cdot\:2^{n-1}\\[-2pt]
\makebox[3zw][c]{$\vdots$}&\hspace{4zw}\vdots\\[-7pt]
\makebox[3zw][c]{$\vdots$}&\hspace{4zw}\vdots
\end{array}\]

\br{0.5}
よって一般項は，$a_n=1\cdot2^{n-1}$と表せる。

\br{0.5}
以上により一般に次のことが成り立つ。

\br{0.5}
\begin{長い枠囲み}{等比数列の一般項}
初項が$a$，公比が$r\ (\neq0)$の等比数列の一般項$\,a_n$は，
\[a_n=a\cdot r^{n-1}\]
\end{長い枠囲み}

すなわち等比数列は初項と公比が決まれば決定される。

\begin{枠囲み}{等比数列}%%%%%%%%%%%%%%%%%%%%%%%%%%%%%%%%%%%%%%%%%%%%%　枠囲み５（J3-2枠囲み８）
第$2$項が$6$，第$5$項が$-48$の等比数列の公比$r$と初項及び第８項を求めよ。\!ただし$r$は実数とする。
\end{枠囲み}
\begin{指針}
2つの与えられた条件から$a$（初項）と$r$に関する連立方程式を導き出す。
\end{指針}

題意の数列の初項を$a$，一般項を$a_n$とすると，
\[a_n=a\cdot r^{n-1}\式番号{1}\]
と表せる。ここで，$a_2=6,\ a_5=-48$より，
\begin{align*}
a_2&=a\cdot r=6\式番号{2}\\
a_5&=a\cdot r^4=-48\式番号{3}
\end{align*}

\ajMaru{3}$\div$\ajMaru{2}より，
\[r^3=-8\]

$r$は実数であるから，
\[r=-2\式番号{4}\]

これを\ajMaru{2}に代入して，
\[a=-3\式番号{5}\]

\ajMaru{1}\ajMaru{4}\ajMaru{5}より，
\[a_n=-3\cdot(-2)^{n-1}\]

よって，
\[a_8=-3\cdot(-2)^7=384\]

以上をまとめて，
\[\text{公比：$\bs{-2}$，初項：$\bs{-3}$，第$8$項：$\bs{384}$}\kotae\]




$a$，$b$，$c$がこの順に等比数列をなすとき，$b$を\textbf{\textgt{等比中項}}といい，$b=ar$，$c=ar^2$より
\[b^2=ac\]
が成り立つ。

\br{0.5}
\begin{枠囲み}{等比中項}
$a$，$b$，$c$が等比数列をなすとき，等比中項$b$は，$b^2=a\cdot c$をみたす。
\end{枠囲み}

\begin{枠囲み}{等差中項・等比中項}%%%%%%%%%%%%%%%%%%%%%%%%%%%%%%%%%%%%%%%　枠囲み６（J3-2枠囲み９）
$1$，$a$，$b$はこの順に等差数列，$1$，$a$，$b^2$はこの順に等比数列をなす。$a\neq b$のとき，$a$，$b$の値を求めよ。
\end{枠囲み}


$1,\ a,\ b$がこの順に等差数列をなすことより，
\[2a=1+b\式番号{1}\]

$1,\ a,\ b^2$がこの順に等比数列をなすことより，
\[a^2=1\cdot b^2\式番号{2}\]

\ajMaru{2}より，
\[a^2-b^2=0\Leftrightarrow (a+b)(a-b)=0\]

$a\neq b$より，$a-b\neq 0$であるから，
\[a+b=0\Leftrightarrow b=-a\式番号{3}\]

\ajMaru{3}を\ajMaru{1}に代入することにより，
\[\bs{a=\bun13,\ b=-\bun13}\kotae\]






\subsection{等比数列の和}%%%%%%%%%%%%%%%%%%%%%%%%%%%%%%%%%%%%%%%%%%%%%%%%%%%%%%%　3-2
初項$a$，公比$r$の等比数列の初項から第$n$項までの和$S_n$は，
\[S_n\:=a+ar+ar^2+\cdots\cdots\cdots+ar^{n-1}\text{　……\ajMaru{1}}\]
と表される。\ajMaru{1}を両辺$r$倍して
\[rS_n=ar+ar^2+\cdots\cdots+ar^{n-1}+ar^n\text{　……\ajMaru{2}}\]

\ajMaru{1}$-$\ajMaru{2}より
\[(1-r)S_n=a-ar^n=a(1-r^n)\text{　……\ajMaru{3}}\]


\bgroup\番号付き箇条書きの空白設定[1zw]{2zw}{1zw}
\番号付け変更[1]{\ajKakkoroman}
\begin{enumerate}
\item $r\neq1$のとき

$r\neq1$なので，\ajMaru{3}の両辺を$1-r$（$\neq0$）で割って
\[S_n=\bunsuu{a(1-r^n)}{1-r}\]

\item $r=1$のとき

\ajMaru{1}は，$S_n=\underbrace{a+a+a+\cdots\cdots\cdots+a}$となるので，

\hspace{10.5zw}$n$個
\[S_n=na\]
\end{enumerate}\egroup%

\br{0.5}
\begin{長い枠囲み}{等比数列の和}
初項$a$，公比$r$の等比数列の初項から第$n$項までの和$S_n$は，

\番号付き箇条書きの空白設定[1zw]{2zw}{1zw}
\番号付け変更[1]{\ajKakkoroman}
\begin{enumerate}
\br{0.4}
\item $r\neq1$のとき　$S_n=\bunsuu{a(1-r^n)}{1-r}$　　$\left[\;\bunsuu{\text{（初項）}\times\{1-\text{（公比）}^{\text{（項数）}}\}}{1-\text{（公比）}}\;\right]$
\br{0.2}
\item $r=1$のとき　$S_n=na$
\end{enumerate}
\end{長い枠囲み}
\ajLig{(注)}　第$n$項\;$ar^{{\scriptsize \uwave{n-1}}}$と第$n$項までの和$\bunsuu{a(1-r^{{\scriptsize \uwave{\,n\;}}})}{1-r}$の$r$の指数の違いに注意。

\begin{枠囲み}{等比数列の和}%%%%%%%%%%%%%%%%%%%%%%%%%%%%%%%%%%%%%%%%%%%%%　枠囲み７（J3-2枠囲み10）
\設問カッコ{1}{$2$，$6$，$18$，……で表される等比数列の初項から第$n$項までの和を求めよ。}
\設問カッコ{2}{初項が$3$，末項が$96$，初項から末項までの和が$189$であるような等比数列の公比と項数を求めよ。ただし，公比は実数であるとする。}
\end{枠囲み}

\begin{指針}
等比数列の和を考えるときは公比$\neq1$を必ず断ること。
\end{指針}



\解答小問{\shoumonKakko{1}%%%%%%%%%%%%%%%%%%%%(1)
題意の数列は初項$2$，公比$3$の等比数列である。

よって，(公比)$\neq 1$に注意して初項から第$n$項までの和$S_n$は，
\[S_n=\bun{2(1-3^n)}{1-3}=\bs{3^n-1}\kotae\]

\shoumonKakko{2}%%%%%%%%%%%%%%%%%%%%(2)
題意の等比数列の公比を$r$，項数を$n$とする。
初項が$3$，末項が$96$であることより，
\[3\cdot r^{n-1}=96\Leftrightarrow r^{n-1}=32\式番号{1}\]

\ajMaru{1}より，$r\neq1$であるから，初項から末項までの和が$\Bun{3(1-r^n)}{1-r}$と表されることより，
\[\bun{3(1-r^n)}{1-r}=189\式番号{2}\]

\ajMaru{1}より，
\[r^n=32r\]

これを\ajMaru{2}に代入して，
\[1-32r=63(1-r)\Leftrightarrow r=2\式番号{3}\]

\ajMaru{1}\ajMaru{3}より，
\[2^{n-1}=32\Leftrightarrow n=6\]

以上より，
\[\text{公比：$\bs{2}$，項数：$\bs{6}$}\kotae\]
}



\subsection{和の公式}%%%%%%%%%%%%%%%%%%%%%%%%%%%%%%%%%%%%%%%%%%%%%%%%%%%%　4-2
\br{0.2}
今までは数列の和を表すのに$\,S_n$という記号を用いてきた。$S_n$はそれなりに便利な記号であるが，一見しただけではどのような項を加えた和なのかはわからない。そこで，加える項をはっきり明示するため，ギリシャ文字$\sum$（シグマ，ローマ字のSに対応する文字で，$\mathrm{Sum}$（和）の頭文字）を導入し，次のように定義する。
\[\sum_{k=1}^{n}a_k=a_1+a_2+\cdots\cdots\cdots+a_n\]

\br{0.5}
ここで，$\sum$の下には加え始める項の番号を，$\sum$の上には加え終わる項の番号を示してある。つまり$\dsum_{k=1}^{n}a_k$は数列$\{\,a_k\,\}$の$k=1$（第1項）から$k=n$（第$n$項）までの和ということである。ゆえに，
\[\sum_{k=1}^{5}a_k=a_1+a_2+a_3+a_4+a_5\hspace{3zw}\sum_{k=3}^{5}a_k=a_3+a_4+a_5\]
のようになる。

\br1
\begin{長い枠囲み}{\ $\bs{\sum}$の性質}

\番号付き箇条書きの空白設定[1zw]{3zw}{1zw}
\begin{originalenumerate}[1.]
\item $\dsum_{k=1}^{n}(a_k+b_k)=\dsum_{k=1}^{n}a_k+\dsum_{k=1}^{n}b_k$

\br{0.3}
\item $\dsum_{k=1}^{n}ca_k=c\dsum_{k=1}^{n}a_k$

\br{0.3}
\item $\dsum_{k=1}^{n}(ca_k+db_k)=c\dsum_{k=1}^{n}a_k+d\dsum_{k=1}^{n}b_k$　（$c$，$d$は定数）

特に，$\dsum_{k=1}^{n}c=\underbrace{c+c+\cdots\cdots\cdots+c}=nc$　　$\dsum_{k=1}^{n}1=n$

\hspace{10.5zw}$n$個
\end{originalenumerate}
\end{長い枠囲み}

%\br1
これは以下のようにして証明される。
\begin{align*}
\hspace{-2zw}1.\quad\makebox[6zw][l]{$\dsum_{k=1}^{n}(a_k+b_k)$}&=(a_1+b_1)+(a_2+b_2)+\cdots\cdots\cdots+(a_n+b_n)\\[-3pt]
&=(a_1+a_2+\cdots\cdots\cdots+a_n)+(b_1+b_2+\cdots\cdots\cdots+b_n)\\
&=\sum_{k=1}^{n}a_k+\sum_{k=1}^{n}b_k
\end{align*}
\begin{align*}
\hspace{-2zw}2.\quad\makebox[6zw][r]{$\dsum_{k=1}^{n}ca_k$\quad}&=ca_1+ca_2+\cdots\cdots\cdots+ca_n\\[-3pt]
&=c(a_1+a_2+\cdots\cdots\cdots+a_n)\\
&=c\sum_{k=1}^{n}a_k
\end{align*}

\begin{長い枠囲み}{和の公式}

\番号付き箇条書きの空白設定[1zw]{3zw}{1zw}
\begin{originalenumerate}[1.]
\item $\dsum_{k=1}^{n}k\hspace{6.5pt}=\;1\;+\:2\;+\:3\;+\cdots\cdots\cdots+\:n\:=\:\bunsuu{1}{2}n(n+1)$

\br{0.3}
\item $\dsum_{k=1}^{n}k^2\,=\;1^2+2^2+3^2+\cdots\cdots\cdots+n^2=\:\bunsuu{1}{6}n(n+1)(2n+1)$

\br{0.3}
\item $\dsum_{k=1}^{n}k^3\,=\;1^3+2^3+3^3+\cdots\cdots\cdots+n^3=\:\left\{\bunsuu{1}{2}n(n+1)\right\}^2$

\br{0.3}
\item $\dsum_{k=1}^{n}r^k\hspace{2pt}=\;r\:+\,r^2+\,r^3+\cdots\cdots\cdots+r^n=\:\bunsuu{r(1-r^n)}{1-r}$　（$r$は$1$でない実数）
\end{originalenumerate}

\end{長い枠囲み}

$\dsum_{k=1}^{n}k$は初項$1$，公差$1$の等差数列の初項から第$n$項までの和であるから，
\br{-0.5}
\[\qquad\sum_{k=1}^{n}k=\bunsuu{1}{2}n(n+1)\]

$\dsum_{k=1}^{n}k^2,\quad\dsum_{k=1}^{n}k^3$についても上の公式を確実に覚えておくこと。

$\dsum_{k=1}^{n}r^k$は初項$r$，公比$r$の等比数列の初項から第$n$項までの和であるから，
\br{-0.5}
\[\qquad\dsum_{k=1}^{n}r^k=\bunsuu{r(1-r^n)}{1-r}\text{　（ただし$\;r\neq1$）}\]




%\newpage
%\br1
%\begin{枠囲み}{和の公式}%%%%%%%%%%%%%%%%%%%%%%%%%%%%%%%%%%%%%%%%%%%%%%%%%%%　枠囲み９（J3-3枠囲み13）
%恒等式$(k+1)^3-k^3=3k^2+3k+1$を利用して，公式$\dsum_{k=1}^{n}k^2=\Bun{1}{6}n(n+1)(2n+1)$を証明せよ。
%\end{枠囲み}
%\begin{指針}
%与えられた恒等式の利用の仕方がポイントである。1つずらしの差の形をしているので，「和の中抜け」に結びつける。
%
%「和の中抜け」とは，一般項$a_n$を$a_n=b_{n+1}-b_n$の形に書き換えて和を求める方法である。このとき，
%\begin{align*}
%\sum_{k=1}^{n}a_k&=\sum_{k=1}^{n}(b_{k+1}-b_k)\\
%&=(b_2-b_1)+(b_3-b_2)+(b_4-b_3)+\cdots\cdots+(b_{n+1}-b_n)\\
%&=b_{n+1}-b_1
%\end{align*}
%となる。本問では，右辺の和に$\dsum_{k=1}^{n}k$，$\dsum_{k=1}^{n}1$が含まれるが，これは計算することができる。
%\end{指針}

\br1
\begin{枠囲み}{和の計算１}%%%%%%%%%%%%%%%%%%%%%%%%%%%%%%%%%%%%%%%%%%%%%%%%%　枠囲み10（J3-3枠囲み14）
次の数列の初項から第$n$項までの和を求めよ。
\二段組小問{\ajKakko{1}　$1^2$，$3^2$，$5^2$，………}{\ajKakko{2}　$1+\Bun{1}{2}$，$2+\bunsuu{1}{4}$，$3+\bunsuu{1}{8}$，$4+\bunsuu{1}{16}$，……}
\end{枠囲み}

\begin{指針}
まず一般項$a_k$を$k$で表し，公式を利用して$\sum{a_k}$を計算する。
\end{指針}

\解答小問{\shoumonKakko{1}%%%%%%%%%%%%%%%%%%%%(1)
与えられた数列の第$k$項は$(2k-1)^2$と表されるので，初項から第$n$項までの和は，
\begin{align*}
&\sum_{k=1}^{n}(2k-1)^2\\
={}&\sum_{k=1}^{n}(4k^2-4k+1)\\
={}&4\sum_{k=1}^{n}k^2-4\sum_{k=1}^{n}k+\sum_{k=1}^{n}1\\
={}&4\cdot\bun16n(n+1)(2n+1)-4\cdot\bun12n(n+1)+n\\
={}&\bs{\bun13n(2n+1)(2n-1)}\kotae
\end{align*}

\shoumonKakko{2}%%%%%%%%%%%%%%%%%%%%(2)
与えられた数列の第$k$項は$k+\bun1{2^k}$と表されるので，初項から第$n$項までの和は，
\begin{align*}
&\sum_{k=1}^{n}\left(k+\bun1{2^k}\right)\\
={}&\sum_{k=1}^{n}k+\sum_{k=1}^{n}\left(\bun12\right)^k\\
={}&\bun12n(n+1)+\bun{\bun12\left\{1-\left(\bun12\right)^n\right\}}{1-\bun12}\\
={}&\bs{\bun12n(n+1)+1-\left(\bun12\right)^n}\kotae
\end{align*}

}


\end{document}